\section{Least-energy source majorants and the Perron--variational gap}
\label{sec:r12-perron-variational-gap}

Retain the full hypotheses (ED1), the original arbitrary background
$S$, and the unchanged strong cone $\mathscr C$. Put $d=q+1$.
For dual functionals write
\[
       \ell\preceq_{\rm psh}\lambda
       \quad\Longleftrightarrow\quad
       \ell(v)\leq\lambda(v)\quad\text{for every }v\in\mathscr C.
 \tag{PV1}
\]
This is the order tested by negative energy potentials. It is
neither ordinary measure order nor middle positivity of currents.

\begin{lemma}[A one-sided Bregman bound]
\label{lem:r12-one-sided-Bregman}
For $a,b\in\mathscr C$,
\[
 \begin{split}
 e(a)-e(b)-dJ(b)(a-b)
  &=\sum_{j=0}^{q-1}(q-j)
       \int_V\Gamma_{\widehat N(a^{[j]},b^{[q-1-j]})}(a-b),\\
 I(a,b)&\leq e(a)-e(b)-dJ(b)(a-b)\leq qI(a,b).
 \end{split}
 \tag{PV2}
\]
Every displayed energy is finite and nonnegative.
\end{lemma}

\noindent\textbf{Status: \texttt{DIRECT\_COROLLARY}} of
the frozen mixed variation and the coefficient computation (JD9).
The bounded Bregman calculation was already present in R3.

\begin{proof}
Apply polynomial Taylor expansion to $e((1-t)b+ta)$ and
retain its first derivative rather than canceling it by
averaging. The same beta integral as in (JD9) gives the
coefficients $q-j$. Their unweighted sum is $I(a,b)$ by
the frozen discrepancy identity, proving both inequalities.
The completed mixed energy calculus justifies all global
integrals before this algebraic computation is made.
\end{proof}

\begin{theorem}[The unique least-energy curvature majorant]
\label{thm:r12-least-energy-curvature-majorant}
Let $\ell\in\mathbb D_{S,\mathrm s}^*$ be nonnegative on
$\mathscr C$, and put $u=\Pi(\ell)$. Then
\[
       \ell\preceq_{\rm psh}J(u),\qquad
       \|J(u)\|_*=\|\ell\|_*.
 \tag{PV3}
\]
Among all $v\in\mathscr C$ satisfying
$\ell\preceq_{\rm psh}J(v)$, the potential $u$ is the
unique one of least energy. Every such $v$ satisfies
\[
               0\leq I(u,v)\leq q\bigl(e(v)-e(u)\bigr).
 \tag{PV4}
\]
Consequently any row of these majorants with
$e(v_j)\to e(u)$ converges strongly to $u$ on the full
original current calculus.

The map $\ell\mapsto J(\Pi(\ell))$ is an idempotent,
positively homogeneous, dual-norm-preserving retraction on
this positive dual cone. Its norm continuity is (NI11).
The phrase \emph{least energy} does not assert a least element
for the order (PV1), and norm preservation does not mean
preservation of the total mass of a measure.
\end{theorem}

\noindent\textbf{Status: \texttt{NEW\_PROVED}} for the
all-degree majorant characterization and its quantitative
energy gap. The normal decomposition is the preceding
proved R12 input.

\begin{proof}
The normal residual in (NI7) is nonpositive on $\mathscr C$,
which proves the order assertion in (PV3). Positivity of
$\ell$ on this cone gives $p(\ell)=\|\ell\|_*$, so (NI9)
and (ED10) prove its norm equality.

For a competitor $v$, the majorant condition tested at $u$
gives $B(u,v^{[q]})=J(v)(u)\geq\ell(u)=e(u)$.
Use (PV2) with $(a,b)=(u,v)$ to obtain
\[
 \begin{split}
 I(u,v)&\leq e(u)-e(v)-d\bigl(B(u,v^{[q]})-e(v)\bigr)\\
       &=e(u)+q e(v)-d B(u,v^{[q]})
                         \leq q(e(v)-e(u)).
 \end{split}
\]
This proves (PV4), minimality, and uniqueness by relative
discrepancy separation. It also proves the approximation
assertion. The retraction, homogeneity and continuity follow
from the corresponding exact properties of (NI7)--(NI11).
No positivity of the signed residual measure is inferred
from its potential-test order.
\end{proof}

For a literal positive Radon datum $\mu$, (PV3) says
\[
      \int_V(-v)\,d\mu\leq\int_V(-v)\,d\widehat\mu_u
               \quad(v\in\mathscr C),\qquad
      \int_V(-u)\,d\mu=e(u).
 \tag{PV5}
\]
One may call this a nonlinear energy sweeping of the datum,
provided (PV1), the specified source class, and the absence
of a total-mass claim remain part of the definition.

\begin{lemma}[The variational source cannot have a positive excess]
\label{lem:r12-variational-positive-excess}
Let $\mu$ be a positive original Radon measure whose literal
functional belongs to $\mathbb D_{S,\mathrm s}^*$, and
put $u=\Pi(\ell_\mu)$. If $\widehat\mu_u\geq\mu$
as original measures, then $\widehat\mu_u=\mu$.
\end{lemma}

\noindent\textbf{Status: \texttt{DIRECT\_COROLLARY}} of
radial variational orthogonality and the frozen completed
weak maximum inequality.

\begin{proof}
For every $a\in\mathscr C$, apply the completed weak
maximum inequality to $a$ and the zero potential, assigning
the equality set to $a$. Since $\max(a,0)=0$, it gives
\[
                   \mathbf1_{\{a=0\}}\widehat\mu_a=0.
 \tag{PV6}
\]
All involved measures also ignore the relative pole set.
Now let $\tau=\widehat\mu_u-\mu\geq0$. It is dominated
by $\widehat\mu_u$, so it ignores both sets just mentioned.
Radial stationarity gives the finite equality
\[
                       \int_V(-u)\,d\tau
                                  =e(u)-\ell_\mu(u)=0.
\]
The integrand is strictly positive and finite on the
remaining full carrier $\{-\infty<u<0\}$. Its increasing
level exhaustion therefore makes $\tau$ vanish there too.
Thus $\tau=0$ as a Radon measure. No subtraction of
infinite unweighted total masses has occurred.
\end{proof}

\begin{theorem}[An exact scalar gap between Perron and variational potentials]
\label{thm:r12-perron-variational-gap}
Let $\mu$ be a positive original Radon measure with a
literal strong-continuous energy functional. Suppose the
completed subsolution family
\[
             \mathscr S_\mu
               =\{v\in\mathscr C:\widehat\mu_v\geq\mu\}
 \tag{PV7}
\]
is nonempty. Let $H_\mu$ be its greatest relative member,
constructed by the frozen R11 full Perron theorem (SE1)--(SE8),
and let $u_\mu=\Pi(\ell_\mu)$. Define
\[
                  \Delta_S(\mu)=e(H_\mu)-e(u_\mu).
 \tag{PV8}
\]
Then
\[
              \Delta_S(\mu)\geq0,\qquad
              I(H_\mu,u_\mu)\leq q\Delta_S(\mu).
 \tag{PV9}
\]
Moreover the following conditions are equivalent:
\begin{enumerate}
\item $\Delta_S(\mu)=0$;
\item $H_\mu=u_\mu$ as relative completed classes;
\item $\widehat\mu_{u_\mu}=\mu$ on every original test;
\item there exists $v\in\mathscr C$ with $\widehat\mu_v=\mu$;
\item there exists a row $v_j\in\mathscr S_\mu$ with
      $e(v_j)\to e(u_\mu)$.
\end{enumerate}
If they hold, this source potential is unique. The family
is nonempty, in particular, whenever $\mu\leq A\mu_\psi$
for an actual first-energy $\psi$ and finite $A\geq0$.

Thus the general subsolution Perron construction and the
general variational construction are linked by a single
nonnegative scalar gap. Its universal vanishing is not
claimed: the compact-source examples (ML22)--(ML28) have
$\Delta_S(\mu)>0$ in every $2\leq q<r$.
\end{theorem}

\noindent\textbf{Status: \texttt{NEW\_PROVED}} for this
all-degree criterion and quantitative comparison. The two
constructions themselves are explicitly frozen or previously
proved inputs. Unqualified vanishing of the gap is
\texttt{FALSE}, not a theorem hidden in the word Perron.

\begin{proof}
The R11 theorem applies to the nonempty family in (PV7),
keeps its measure lower bound, and makes its relative
envelope an actual member of $\mathscr C$. Thus $H_\mu$
is the greatest member of that family. Every measure
majorant is also a majorant for the potential-test order
(PV1). Applying (PV4) proves (PV9).

If the gap is zero, (PV9) gives equality of the two
relative potentials. The common potential is a subsolution,
so Lemma~\ref{lem:r12-variational-positive-excess} makes
its source exactly $\mu$. Conversely, a source potential
$v$ has zero normal residual and hence is the unique
variational potential $u_\mu$. It belongs to (PV7), so
$H_\mu\succeq u_\mu$. Energy monotonicity gives
$e(H_\mu)\leq e(u_\mu)$, and (PV9) gives equality.
This proves the first four equivalences and uniqueness.

Every $v\in\mathscr S_\mu$ is below $H_\mu$, so
$e(v)\geq e(H_\mu)$. A row in condition (5) therefore
forces the gap to be zero. If the gap is zero, the
constant row $v_j=u_\mu$ supplies condition (5).
Alternatively (PV4) directly shows that any such row
converges strongly to $u_\mu$ and its source inequalities
pass to all original compact tests.

For the printed dominated datum, $A^{1/q}\psi$ is a
subsolution (the zero case is immediate), and its literal
datum functional satisfies the strong-dual hypothesis by
mixed Holder and (EV1). The nonlinear compact-source
examples satisfy these same assumptions but have no source
in $\mathscr C$. The equivalence therefore makes their
gap strictly positive. None of these arguments needs a
common actual lower barrier for the whole subsolution family.
\end{proof}

\begin{proposition}[A quantitative positive-test lower bound for the gap]
\label{prop:r12-positive-test-gap-bound}
Under the preceding theorem put $u=u_\mu$, $H=H_\mu$,
$R=r(H)$, and define
\[
 \varepsilon_+(\mu,u)=
   \sup_{\substack{0\leq\phi\in A_c^{0,0}(V)\\
                   \|\phi\|_{\rm ed}\leq1}}
                   \left(\int_V\phi\,d(\mu-\widehat\mu_u)\right)_+.
 \tag{PV10}
\]
This is finite. If $R>0$, then
\[
       \Delta_S(\mu)\geq
          \frac{\varepsilon_+(\mu,u)^2}{q^2R^{q-1}}.
 \tag{PV11}
\]
If $R=0$, the datum and all quantities are zero. Also
$\varepsilon_+(\mu,u)=0$ if and only if the source equation
holds. Every positive test in (PV10) is from the full
original sheaf, not a selected moment family.
\end{proposition}

\noindent\textbf{Status: \texttt{DIRECT\_COROLLARY}} of
the new gap theorem and the proved source-dual estimate.

\begin{proof}
All original compact smooth scalars lie in $\mathbb D_S$
by (HC14)--(HC15). Finiteness follows by the dual norm of
$\ell_\mu-J(u)$. Since $\widehat\mu_H\geq\mu$, each
nonnegative test in (PV10) gives
\[
 \int\phi\,d(\mu-\widehat\mu_u)
       \leq\int\phi\,d(\widehat\mu_H-\widehat\mu_u)
       =(J(u)-J(H))(\phi).
\]
Consequently $\varepsilon_+\leq\|J(H)-J(u)\|_*$.
The energy minimum gives $r(u)\leq R$. Equations (NI2)
and (PV9) now give
$\varepsilon_+\leq qR^{(q-1)/2}\Delta_S(\mu)^{1/2}$,
which is (PV11). When $R=0$, the subsolution $H=0$
already forces $\mu=0$.

Vanishing in (PV10), together with original positive
scalar test detection, says $\widehat\mu_u\geq\mu$.
Lemma~\ref{lem:r12-variational-positive-excess} upgrades
this to equality. The reverse implication is immediate.
This is a full original-measure statement, not an
ordinary-to-delta-test promotion.
\end{proof}

